\section{Introduction}\label{sec:intro}

Grid carbon intensity (CI) varies within a day, so a GPU job that can wait a few hours can often
run on cleaner electricity. Temporal and spatial shifting are now standard
proposals for greener data centres and ML clusters~\cite{wiesner2021wait,radovanovic2023carbon,xu2025green,hanafy2024gaia}.
Evaluations usually report aggregate savings. They rarely ask what happens to the individual job
that was told to wait.

We measure that directly. A deferred job is \emph{harmed} if it ends up emitting more than if it had
simply run on arrival. The common intuition is that harm comes from forecast error: the scheduler
expected a cleaner hour that never arrived. The natural remedy is uncertainty quantification, such as
conformal intervals~\cite{angelopoulos2023gentle,gibbs2021aci,li2024uncertainty}, so that the scheduler
defers only when a saving is likely. We started from that hypothesis and built such schedulers. The
data point elsewhere.

\textbf{Findings.} Replaying ten weeks of the Alibaba PAI GPU trace~\cite{weng2022mlaas} against six
grids, with forecasts that include NESO's own published 48\,h forecasts:
\begin{itemize}
\item Harm is common. A point-forecast scheduler harms \res{V2U}{ALL-S12-utiltarget0.7}{P3}{harmrate}\%
of shifted jobs at 70\% utilisation. A clairvoyant greedy scheduler (CG), which knows future CI
exactly but places jobs one at a time, still harms \res{V2U}{ALL-S12-utiltarget0.7}{P5}{harmrate}\%.
\item Harm is mainly contention. CG's harm rate rises from
\res{V2U}{ALL-S12-utiltarget0.5}{P5}{harmrate}\% to \res{V2U}{ALL-S12-utiltarget0.9}{P5}{harmrate}\%
between 50\% and 90\% utilisation. An exact decomposition (Eq.~\ref{eq:decomp}) attributes
\res{V2U}{POOL-S12-utiltarget0.7}{P2}{sharecontention}\% of a greedy scheduler's harmful excess to jobs
not getting the slot they wanted, and \res{V2U}{POOL-S12-utiltarget0.7}{P2}{shareforecast}\% to forecast error.
\item Uncertainty-aware rules are a dial, not a fix. Conformal policies move along a savings--harm
curve. The decision-level rule's calibration bounds the \emph{marginal} rate of harmful deferrals
over all decisions, and this bound holds in every region at every $\alpha$. It does not bound the
\emph{conditional} harm rate among deferred jobs, which exceeds $\alpha$ in several regions, most
in MISO and PJM.
\item Addressing contention costs less saving per unit of harm removed. Keeping a quarter of capacity free for
later arrivals is not dominated by any point of the conformal dial in \res{V2F}{ALL}{P8}{nondominated} of
\res{V2F}{ALL}{P8}{nregions} regions.
\end{itemize}

\textbf{Contributions.} (i)~A job-level harm metric with an exact three-way attribution to forecast
error, re-planning and contention. (ii)~A measurement across utilisation, slack, forecast error and
region, anchored on operator forecast errors that we did not construct. (iii)~An evaluation of
conformal risk-aware deferral as a savings--harm dial, with its guarantee scoped precisely.
(iv)~A simple congestion-aware rule, which is our constructive result, and guidance on when shifting
is worth it.
